\begin{abstract}
Geometric spatial reasoning underlies applications from robotics to automated design, yet the ability of large language models (LLMs) to turn a verbal description into a physically connected three-dimensional assembly remains largely underexplored. Brick assembly is a demanding constructive probe of this ability: a short, underspecified prompt leaves dimensions, proportions, and part layout implicit, and a valid answer must match the semantic intent of the prompt while completing this geometry with discrete pieces that interlock through studs. We study this setting as the \texttt{Text-to-Brick} task and present BrickForma, a framework that crystallizes the geometric knowledge of large language models into a compact model for this task. Parametric programs written with large language model assistance generate 80,000 source voxel models across 21 categories on a $48^3$ grid, each described at five levels of detail, yielding approximately 400,000 caption records that tie varying amounts of language to shared geometry. A voxel-to-brick conversion pipeline with connection analysis and local refinement turns these shapes into geometrically checked brick targets, on which we fine-tune Qwen3-1.7B with low-rank adaptation to predict brick dimensions and positions from text. A mixed-integer linear programming stage then seeks an exact cover of the voxel target reconstructed from these predictions with a single stud-connected component. On 525 shared prompts spanning five description levels, BrickForma's raw outputs achieve the highest mean VQAScore, PE, and SigLIP2 at every level among the three systems. BrickForma has a mean invalid-placement ratio of 0.79\%, compared with 0.00\% for BrickGPT and 18.42\% for BrickNet. These results suggest that executable programs and checked geometric data are an effective route for carrying spatial knowledge from large models into small models without the need of human-verified input for training.
\end{abstract}